% Image credits for the photographs and AI illustrations of Book 3 (University, Year 1).
% Machine-readable license records live in images/book3/CREDITS.md; keep the
% two files in sync. The Book 3 writer owns this file: add one \item per
% photograph (inside the itemize, by chapter) as photographs land.
\chapter*{Créditos de las imágenes}

Los dibujos, mapas y gráficos de este libro son originales; los mapas y gráficos
elaborados a partir de datos publicados citan su fuente en el pie de figura o en
esta lista. Las fotografías se usan bajo sus respectivas licencias libres, con
gratitud a sus autores; los enlaces completos a la fuente y las URL de las
licencias de cada fotografía figuran en \texttt{images/book3/CREDITS.md} en el
repositorio del libro.

\bigskip
Junto a las fotografías y a los dibujos originales, algunas figuras son
ilustraciones generadas por IA, producidas con la generación de imágenes de OpenAI
a partir de instrucciones escritas por los editores del libro, y revisadas por su
exactitud científica antes de incluirlas. Muestran solo escenas genéricas, nunca un
lugar real, una persona real ni un ejemplar real. La instrucción completa de cada
imagen generada figura en \texttt{images/book3/ai/PROMPTS.md} en el repositorio.

\bigskip
\noindent Fotografías y datos, por capítulo:
\begin{itemize}\small
  \item Cap.~1 (astronomía esférica): estelas de estrellas sobre el Observatorio de Paranal, R.~Wesson/ESO, CC BY 4.0.
  \item Cap.~2 (telescopios, detectores y fotometría): seis segmentos del espejo primario del telescopio espacial James Webb antes de las pruebas criogénicas, NASA/MSFC/David Higginbotham, dominio público.
  \item Cap.~3 (planetología comparada): los planetas terrestres a escala, NASA; Venus a partir del radar de Magallanes, NASA/JPL; Marte desde el orbitador Viking 1, NASA/USGS; Europa, NASA/JPL/DLR; Io, NASA/JPL/University of Arizona; todos de dominio público.
  \item Cap.~4 (mapas geológicos): William Smith, \emph{A Delineation of the Strata of England and Wales, with part of Scotland}, 1815; dominio público (Wikimedia Commons).
  \item Cap.~5 (cortes geológicos): la discordancia de Hutton en Siccar Point, PawełS, CC BY-SA 4.0 (Wikimedia Commons).
  \item Cap.~6 (métodos de campo y de laboratorio): un granito de biotita en lámina delgada entre polarizadores cruzados, Darla Sondrol, CC0 (Wikimedia Commons).
  \item Cap.~7 (minerales formadores de rocas): ejemplares de olivino, cuarzo, biotita, almandino y hornblenda, Robert M.~Lavinsky, CC BY-SA 3.0; labradorita, Twyla Baker, CC0; romboedros de exfoliación de calcita, Giovanna Canu y Eva Santini, CC BY 4.0 (todos en Wikimedia Commons).
  \item Cap.~8 (mineralogía óptica): un gabro en lámina delgada con un polarizador y entre polarizadores cruzados, Chmee2, CC BY-SA 3.0; Auguste Michel-Lévy de estudiante, dominio público (Wikimedia Commons). Los colores de interferencia se calculan a partir de las funciones de correspondencia de color CIE 1931 y del iluminante patrón CIE D65, \copyright{} CIE, CC BY-SA 4.0.
  \item Cap.~9 (la Tierra estratificada): Andrija Mohorovičić, dominio público; Inge Lehmann en 1932, fotografía de Even Neuhaus, dominio público (Wikimedia Commons). PREM: A.~M.~Dziewonski y D.~L.~Anderson (1981), tal como lo tabula ObsPy.
  \item Cap.~10 (el balance térmico de la Tierra): flujo de calor global de J.~H.~Davies y D.~R.~Davies (2010), \emph{Solid Earth}, CC BY 3.0; adiabática del manto de T.~Katsura (2022), \emph{JGR Solid Earth}, CC BY.
  \item Cap.~11 (isostasia y anomalías gravimétricas): el aparato de péndulo de Vening Meinesz, Broodt, CC BY-SA 4.0 (Wikimedia Commons). Altitudes de ETOPO1 (NOAA National Centers for Environmental Information), muestreadas cada grado.
  \item Cap.~12 (cinemática de placas sobre una esfera): una estación permanente de posicionamiento por satélite, Alpsdake, CC BY-SA 4.0 (Wikimedia Commons). Límites de placas de PB2002 (P.~Bird, 2003; GeoJSON por H.~Ahlenius, ODC-By 1.0); líneas de costa de Natural Earth (dominio público).
  \item Cap.~13 (rifts, dorsales y márgenes pasivos): fisuras abiertas en Þingvellir, Islandia, Olga Ernst, CC BY-SA 4.0; lavas almohadilladas en la dorsal de Galápagos, NOAA Okeanos Explorer Program, CC BY 2.0 (ambas en Wikimedia Commons). Profundidades de ETOPO1 (NOAA National Centers for Environmental Information) y edades del fondo oceánico de M.~Seton y otros (2020), EarthByte, CC BY 4.0.
  \item Cap.~14 (márgenes activos): erupción del Sarychev Peak, vista desde la Estación Espacial Internacional, NASA, dominio público (Wikimedia Commons). Terremotos del USGS ANSS Comprehensive Catalog (dominio público); volcanes del Smithsonian Institution Global Volcanism Program; límites de placas de PB2002 (P.~Bird, 2003; GeoJSON por H.~Ahlenius, ODC-By 1.0); costas de Natural Earth (dominio público); edades del fondo oceánico de M.~Seton y otros (2020), EarthByte, CC BY 4.0.
  \item Cap.~15 (puntos calientes y volcanismo intraplaca): una fuente de lava en Pu\textquoteleft u \textquoteleft O\textquoteleft o, Kilauea, G.~E.~Ulrich, U.S. Geological Survey, dominio público; los Ghats Occidentales en Matheran, Nicholas (Nichalp), CC BY-SA 2.5 (ambas en Wikimedia Commons). Posiciones de los puntos calientes, acimutes de las trazas y tasas de la tabla~1 del suplemento electrónico de W.~J.~Morgan y J.~P.~Morgan (2007); profundidades de ETOPO1 (NOAA National Centers for Environmental Information); límites de placas de PB2002 (P.~Bird, 2003; GeoJSON por H.~Ahlenius, ODC-By 1.0); costas de Natural Earth (dominio público).
  \item Cap.~16 (esfuerzo, deformación y reología): una celda triaxial de alta presión para medir la resistencia de las rocas, Darryl Peroni, CSIRO, CC BY 3.0 (Wikimedia Commons). Parámetros de fluencia por dislocación de cuarcita húmeda (Gleason y Tullis, 1995) y de olivino seco (Hirth y Kohlstedt, 2003), tal como los recopila M.~A.~Lopez-Sanchez en \emph{strength\_envelopes} (MPL-2.0).
  \item Cap.~17 (fallas, pliegues y texturas): un espejo de falla en Corona Heights, San Francisco, Dextersinister, CC BY-SA 4.0; brecha de falla del cabalgamiento de Keystone, Nevada, Jstuby, dominio público; pliegues en chevron en Millook, Cornualles, Smalljim, CC BY 3.0; un anticlinal con clivaje axial-plano cerca de Potes, Cantabria, Mikenorton, CC BY-SA 3.0 (todos en Wikimedia Commons).
  \item Cap.~18 (la fuente del terremoto): terreno y vallas separados por la ruptura de 1906 en el condado de Marin, \emph{Popular Science Monthly} vol.~69 (1906), dominio público; desplazamiento del arroyo Wallace a través de la falla de San Andrés, R.~E.~Wallace, U.S. Geological Survey, dominio público (ambas en Wikimedia Commons). Terremotos del USGS ANSS Comprehensive Catalog (dominio público); el mecanismo focal del terremoto de Tohoku de 2011 a partir del tensor momento de fase W del USGS (dominio público).
  \item Cap.~19 (meteorización y arcillas): granito descompuesto y arena granítica en Targassonne, Pirineos, Pinpin, CC~BY~2.5; bauxita sobre arenisca caolinítica en Weipa, Australia, Werner Schellmann, CC~BY-SA~2.5; cristales de caolinita al microscopio electrónico de barrido, U.S. Geological Survey, dominio público (todos en Wikimedia Commons). Radios iónicos de R.~D.~Shannon (1976).
  \item Cap.~20 (erosión y transporte): la pluma de inundación del Misisipi desde el instrumento MODIS del satélite Terra, 4~de marzo de 2018, imagen del NASA Earth Observatory de Joshua Stevens, dominio público (Wikimedia Commons).
  \item Cap.~21 (sedimentos clásticos y sus ambientes): estratificación cruzada en la arenisca de Aztec, Nevada, Chip Carroon, Bureau of Land Management, CC~BY-SA~2.0; el delta del río Lena desde Landsat~7, U.S. Geological Survey, dominio público; una turbidita con las divisiones de la secuencia de Bouma, Becke-Oese, Jo Weber, CC~BY-SA~3.0 (todas en Wikimedia Commons).
  \item Cap.~22 (sedimentos carbonatados, silíceos y evaporíticos): Bora Bora desde la Estación Espacial Internacional, NASA Johnson Space Center, dominio público; una célula de \emph{Emiliania huxleyi}, Alison R.~Taylor, University of North Carolina Wilmington Microscopy Facility, CC~BY~2.5; \emph{Orbulina universa}, H.~Takagi, K.~Kimoto, T.~Fujiki, H.~Saito, C.~Schmidt, M.~Kucera y K.~Moriya, CC~BY~4.0; diatomeas del estrecho de McMurdo, Gordon T.~Taylor, Stony Brook University, NOAA Corps Collection, dominio público (todas en Wikimedia Commons).
  \item Cap.~23 (diagénesis y cuencas sedimentarias): crecimientos de cuarzo en la arenisca de Potsdam, M.~C.~Rygel, CC~BY-SA~4.0 (Wikimedia Commons). Apilamiento del nivel del mar de R.~M.~Spratt y L.~E.~Lisiecki (2016), NOAA NCEI World Data Service for Paleoclimatology; parámetros de compactación de J.~G.~Sclater y P.~A.~F.~Christie (1980), tal como los tabula pyBacktrack (EarthByte).
  \item Cap.~24 (la escala estratigráfica): clavo dorado del GSSP del Selandiense en Zumaia, Josugoni, CC~BY-SA~4.0 (Wikimedia Commons); el límite Cretácico--Paleógeno en el Parque Estatal Trinidad Lake, Nationalparks, CC~BY-SA~2.5 (Wikimedia Commons). Edades base y colores de la Carta Cronoestratigráfica Internacional, Comisión Internacional de Estratigrafía, CC~BY~4.0.
  \item Cap.~25 (métodos de datación radiométrica): un grano de circón en catodoluminiscencia, Emmanuel Roquette, CC~BY-SA~4.0; un espectrómetro de masas de ionización térmica, Radiogenic, CC0 (ambas en Wikimedia Commons). La curva de calibración IntCal20 de P.~J.~Reimer y otros (2020), \emph{Radiocarbon} 62, CC~BY~4.0, tal como se distribuye con el paquete rintcal (M.~Blaauw).
  \item Cap.~26 (estructura de la atmósfera y del océano): recuperación de un CTD y una roseta en el estrecho de Bering, NASA Goddard Space Flight Center, dominio público (Wikimedia Commons). Temperatura y salinidad del World Ocean Atlas 2023 (NOAA National Centers for Environmental Information).
  \item Cap.~27 (radiación y transporte de calor): la Tierra iluminada desde el Deep Space Climate Observatory, NASA Goddard Space Flight Center, dominio público (Wikimedia Commons). Radiación, flujos superficiales y precipitación del NCEP-NCAR Reanalysis~1 (E.~Kalnay y otros, 1996), datos proporcionados por NOAA PSL, Boulder, Colorado, EE.~UU., desde su sitio web en https://psl.noaa.gov.
  \item Cap.~28 (la circulación general de la atmósfera): un ciclón extratropical en el golfo de Alaska, NASA MODIS/LANCE, dominio público; la zona de convergencia intertropical a partir de datos de nubes del GOES-11, NASA Earth Observatory, dominio público (ambas en Wikimedia Commons). Vientos y presión a nivel del mar del NCEP-NCAR Reanalysis~1 (E.~Kalnay y otros, 1996), datos proporcionados por NOAA PSL, Boulder, Colorado, EE.~UU., desde su sitio web en https://psl.noaa.gov.
  \item Cap.~29 (la circulación del océano): la corriente del Golfo vista a través de la temperatura de la superficie del mar, NOAA NESDIS, dominio público; un flotador perfilador de océano profundo en despliegue, fotógrafo no identificado de NOAA, dominio público (ambas en Wikimedia Commons).
\end{itemize}
